\section{Metodología}\label{metodologuxeda}

El desarrollo del proyecto se gestionará utilizando las metodologías Scrum y LEAN. Scrum, que de hecho está basado en la filosofía LEAN, permitirá organizar y dar seguimiento a las actividades necesarias para construir el producto mínimo viable, además de permitir la mejora continua tanto del proceso de desarrollo como del equipo de trabajo, mientras que LEAN proporcionará un marco de validación orientado a verificar que la solución propuesta genera valor para la organización. La aplicación conjunta de ambas metodologías permitirá gestionar el desarrollo de manera iterativa, medir resultados mediante indicadores cuantificables y reducir el riesgo asociado a la implementación de la propuesta de innovación.

\subsection{Roles, Artefactos y Eventos}\label{roles-artefactos-y-eventos}

Scrum establece que se debe definir el rol que tendrá cada miembro del equipo, los artefactos y herramientas que se utilizarán y los eventos que permitirán revisar y hacer los cambios necesarios para garantizar la calidad del producto. Estos elementos, en el marco del proyecto, se describen a continuación.

\subsubsection{Roles}\label{roles}

La metodología define tres roles \citep{schwaber2020}: el propietario del producto (Product Owner), responsable de las decisiones que permiten lograr el máximo valor para el producto; el Scrum Master, encargado de dirigir las reuniones del equipo y de ayudar a sus miembros a mejorar su flujo de trabajo; y el equipo de desarrollo, que en cada iteración implementa los incrementos hasta completar el producto final. Tomando en cuenta que el equipo de trabajo es pequeño, formado por cuatro integrantes, se han definido los siguientes roles:

\begin{itemize}
\item
  \textbf{Product Owner}: Francisco Camacho, considerando que la propuesta beneficiará a una empresa de Bolivia y es el miembro del equipo que está más cerca y está en contacto directo y permanente con la empresa, su gerente y sus operadores.
\item
  \textbf{Scrum Master}: Mario Diaz, elegido por el equipo por sus habilidades de gestión y dirección.
\item
  \textbf{Equipo de Desarrollo}: Los cuatro miembros del equipo: Francisco Camacho, Mario Diaz, Angel Hernández y Guillermo Sais de la Torre. En el caso de Francisco y Mario, que cumplirán un doble rol, dedicarán el 75 \% de su tiempo a las tareas de desarrollo y 25 \% de su tiempo a las tareas de gestión en el rol asignado.
\end{itemize}

La Tabla~\ref{tab:roles-del-equipo-del-trabajo} resume los roles y responsabilidades del equipo de trabajo:

\captionof{table}{Roles del equipo del trabajo en el proceso de desarrollo}\label{tab:roles-del-equipo-del-trabajo}

\begin{longtable}[]{@{}
  >{\raggedright\arraybackslash}p{(\columnwidth - 4\tabcolsep) * \real{0.3268}}
  >{\raggedright\arraybackslash}p{(\columnwidth - 4\tabcolsep) * \real{0.2483}}
  >{\raggedright\arraybackslash}p{(\columnwidth - 4\tabcolsep) * \real{0.4248}}@{}}
\toprule\noalign{}
\begin{minipage}[b]{\linewidth}\raggedright
\textbf{Integrante}
\end{minipage} & \begin{minipage}[b]{\linewidth}\raggedright
\textbf{Rol}
\end{minipage} & \begin{minipage}[b]{\linewidth}\raggedright
\textbf{Responsabilidades}
\end{minipage} \\
\midrule\noalign{}
\endhead
\bottomrule\noalign{}
\endlastfoot
Francisco J. Camacho Dorado & Product Owner y Equipo de Desarrollo. & Definir prioridades; validar entregables; asegurar el alineamiento con las necesidades de Cafetalino; desarrollo de modelos analíticos; documentación técnica. \\
Mario Alberto Díaz Zetina & Scrum Master y Equipo de Desarrollo. & Coordinar las actividades del equipo; facilitar reuniones; gestionar impedimentos; investigación; análisis de datos; validación de resultados. \\
Ángel Zait Hernández López & Equipo de Desarrollo. & Desarrollo de modelos analíticos; documentación técnica. \\
Guillermo Sais de la Torre & Equipo de Desarrollo. & Investigación; análisis de datos; validación de resultados. \\
\end{longtable}

\fuente{Fuente: Elaboración propia}

\subsubsection{Artefactos}\label{artefactos}

Scrum especifica tres artefactos principales \citep{schwaber2020}: la lista de tareas pendientes del proyecto (Product Backlog), que ordena por prioridad todo lo que se requiere para construir el producto; el registro de tareas de una iteración (Sprint Backlog), que es el plan de trabajo de ese sprint; y el incremento, una versión parcial del producto que ya funciona. Los tres se describen a continuación en el marco del proyecto.

\textbf{Lista de tareas pendientes del proyecto (Product Backlog)}

La lista de tareas pendientes del proyecto se gestionará a través de Jira, como muestra la Figura \ref{fig:lista-de-tareas-pendientes-del} del Anexo~\ref{anexo-capturas-de-jira}, y reúne las historias de usuario necesarias para desarrollar el sistema propuesto para Cafetalino, cuyos componentes son: el análisis exploratorio de los datos disponibles, el modelo predictivo de reabastecimiento, el modelo de enrutamiento dinámico y la interfaz de usuario para interactuar con ambos. Cada historia se refina en elementos más pequeños y precisos, se describe, se le asigna un orden de prioridad y se estima su nivel de complejidad antes de poder incorporarse a una iteración \citep{schwaber2020}.

La lista inicial completa se recoge en la Tabla~\ref{tab:backlog-inicial-prioridad-y-nivel} del Anexo~\ref{anexo-backlog-inicial}, con las historias agrupadas por épica, sus criterios de aceptación, su prioridad y la estimación de su nivel de complejidad. Esta última se expresa en puntos de historia (Story Points, SP) siguiendo la serie de Fibonacci, un criterio que obliga a diferenciar tamaños relativos en lugar de intentar precisiones ilusorias; la asignación se realiza en consenso con todo el equipo y pondera la dificultad técnica de la tarea, la incertidumbre asociada a ella y el esfuerzo estimado. La estimación sitúa el techo de dificultad en dos historias de 8 puntos —la construcción del modelo predictivo y el cálculo de la ruta más corta—, que son las que concentran la incertidumbre técnica del proyecto; el resto se reparte entre 2 y 5 puntos, con los valores más bajos en las historias de presentación, que se apoyan en componentes ya conocidos por el equipo.

\textbf{Sprint Backlog}

El registro de tareas de cada iteración —el plan que los desarrolladores elaboran para guiar su trabajo, con el objetivo del sprint y el subconjunto de historias seleccionadas \citep{schwaber2020}— se extraerá de esa misma lista y se manejará también en Jira, tal como se aprecia en la Figura \ref{fig:registro-de-tareas-de-una}, donde las historias aparecen clasificadas según su estado: pendiente, en ejecución, en revisión y terminada. La herramienta permitirá garantizar que cada historia de usuario se ha refinado, se le ha asignado un orden de prioridad, se ha estimado su nivel de complejidad y, fundamentalmente, que cumple las definiciones de Listo del apartado \ref{definiciones-de-listocompletado-readydone}. Cada iteración se dedicará a uno de los cuatro componentes ya enunciados, en ese mismo orden: el análisis exploratorio del historial de recargas, el modelo predictivo a partir de datos históricos y sin hardware adicional, el modelo de enrutamiento dinámico construido sobre la demanda proyectada y la interfaz web que da acceso a ambos.

\textbf{Incremento}

Cada incremento es una versión parcial del producto que ya funciona, es aditivo a los anteriores y debe verificarse a fondo para garantizar que todos operen en conjunto \citep{schwaber2020}. En el sistema propuesto, cada uno de los cuatro componentes anteriores será un incremento, que deberá cumplir las definiciones de completado del apartado \ref{definiciones-de-listocompletado-readydone} y ser integrado y probado junto con los incrementos ya desarrollados y verificados.

\subsubsection{Eventos}\label{eventos}

Scrum contempla cinco eventos \citep{schwaber2020}: la iteración o Sprint, que hace de contenedor para los otros cuatro; la reunión de planificación de la iteración (Sprint Planning); la reunión diaria de Scrum (Daily Scrum), de quince minutos, para revisar avances, resolver atascos y ajustar el plan; la reunión de revisión de la iteración (Sprint Review), donde se examina el resultado y se definen los ajustes necesarios para alcanzar el producto final; y la retrospectiva (Sprint Retrospective), dedicada al desempeño del propio equipo. Todas las reuniones del proyecto se realizarán por Google Meet en horario nocturno, para no interrumpir el trabajo cotidiano de los integrantes, y se mantendrá una comunicación permanente de coordinación por WhatsApp.

Para desarrollar el sistema de reabastecimiento predictivo y enrutamiento dinámico se ha establecido que cada iteración dure dos semanas y se han proyectado cuatro iteraciones, es decir ocho semanas efectivas de trabajo entre agosto y septiembre de 2026, tal como se observa en la Figura \ref{fig:iteraciones-del-proyecto} del Anexo~\ref{anexo-capturas-de-jira}. En una iteración no se harán cambios que pongan en riesgo su objetivo o disminuyan la calidad, aunque el alcance podrá aclararse o renegociarse con el Product Owner conforme se aprenda más sobre el producto \citep{schwaber2020}.

Las cuatro reuniones serán dirigidas por Mario Diaz (Scrum Master) y atenderán los lineamientos de Francisco Camacho (Product Owner). Las historias comprometidas en la planificación se subdividirán en piezas implementables en un día como máximo, y la retrospectiva se celebrará antes de la planificación siguiente, de modo que sus acuerdos se apliquen ya en la iteración que se abre. La Tabla~\ref{tab:eventos-de-scrum-para-el} resume la frecuencia, la duración y el objetivo de cada evento.

\captionof{table}{Eventos de Scrum para el Proyecto}\label{tab:eventos-de-scrum-para-el}

\begin{longtable}[]{@{}
  >{\raggedright\arraybackslash}p{(\columnwidth - 6\tabcolsep) * \real{0.1720}}
  >{\raggedright\arraybackslash}p{(\columnwidth - 6\tabcolsep) * \real{0.2420}}
  >{\raggedright\arraybackslash}p{(\columnwidth - 6\tabcolsep) * \real{0.1401}}
  >{\raggedright\arraybackslash}p{(\columnwidth - 6\tabcolsep) * \real{0.4459}}@{}}
\toprule\noalign{}
\begin{minipage}[b]{\linewidth}\raggedright
\textbf{Evento}
\end{minipage} & \begin{minipage}[b]{\linewidth}\raggedright
\textbf{Frecuencia / Agenda}
\end{minipage} & \begin{minipage}[b]{\linewidth}\raggedright
\textbf{Duración estimada}
\end{minipage} & \begin{minipage}[b]{\linewidth}\raggedright
\textbf{Objetivo}
\end{minipage} \\
\midrule\noalign{}
\endhead
\bottomrule\noalign{}
\endlastfoot
\textbf{Sprint Planning} & Inicio de cada sprint.

Lunes 19:00 - 20:00 & 60 min & Definir las historias de usuario a desarrollar en el sprint, revisar prioridades, estimar esfuerzo y asignar responsabilidades. \\
\textbf{Daily Scrum} & Dos veces / semana.

Lunes y jueves

18:30 - 18:45 & 15 min & Revisar el avance, dificultades y coordinar el trabajo del equipo. \\
\textbf{Sprint Review} & Final de cada sprint.

Sabado 18:00 - 18:30 & 30 min & Presentar los resultados, verificar el cumplimiento de las historias de usuario comprometidas, verificar indicadores de éxito del proyecto y validar el alineamiento con los objetivos del proyecto. \\
\textbf{Sprint Retrospective} & Final de cada sprint.

Lunes 18:30 - 19:00 & 30 min & Identificar oportunidades de mejora y definir acciones correctivas para el desempeño del equipo. \\
\end{longtable}

\fuente{Fuente: Elaboración propia}

\subsection{Definiciones de Listo/Completado (Ready/Done)}\label{definiciones-de-listocompletado-readydone}

La metodología Scrum obliga a declarar explícitamente la Definición de Listo (Definition of Ready, DoR) y la Definición de Completado (Definition of Done, DoD): los criterios que establecen cuándo una historia de usuario puede empezar a desarrollarse y cuándo está terminada. Ambas son fundamentales para garantizar la calidad del software y minimizar las historias incompletas o mal implementadas, que obligan a repetir el trabajo \citep{schwaber2020,plaskonka2023}. En este apartado se formalizan las dos definiciones que se utilizarán en el proyecto.

\subsubsection{Definición de Listo (Definition of Ready - DoR)}\label{definiciuxf3n-de-listo-definition-of-ready---dor}

La Definición de Listo reúne los requisitos que una historia de usuario debe cumplir para pasar de la lista de tareas pendientes del proyecto al registro de tareas de una iteración \citep{plaskonka2023}. En el presente proyecto, una historia será considerada lista para incorporarse a una iteración cuando cumpla las condiciones que enumera la Tabla~\ref{tab:definicion-de-listo-dor-para}:

\captionof{table}{Definición de Listo (DoR) para el proyecto}\label{tab:definicion-de-listo-dor-para}

\begin{longtable}[]{@{}
  >{\raggedright\arraybackslash}p{(\columnwidth - 2\tabcolsep) * \real{0.3141}}
  >{\raggedright\arraybackslash}p{(\columnwidth - 2\tabcolsep) * \real{0.6859}}@{}}
\toprule\noalign{}
\begin{minipage}[b]{\linewidth}\raggedright
\textbf{Criterio}
\end{minipage} & \begin{minipage}[b]{\linewidth}\raggedright
\textbf{Descripción}
\end{minipage} \\
\midrule\noalign{}
\endhead
\bottomrule\noalign{}
\endlastfoot
\textbf{Objetivo claramente definido} & La funcionalidad o entregable requerido se encuentra descrito y comprendido por el equipo. \\
\textbf{Criterios de aceptación establecidos} & Existen condiciones verificables que permitan determinar el éxito de la historia de usuario. \\
\textbf{Fuente de información identificada} & Los datos, documentos o recursos necesarios se encuentran disponibles. \\
\textbf{Dependencias identificadas} & Se conocen las actividades o entregables previos necesarios para iniciar el trabajo. \\
\textbf{Esfuerzo estimado} & La historia cuenta con una estimación acordada por el equipo. \\
\textbf{Valor para el proyecto} & La actividad contribuye directamente al cumplimiento de los objetivos del MVP. \\
\end{longtable}

\fuente{Fuente: Elaboración propia}

\subsubsection{Definición de Completado (Definition of Done - DoD)}\label{definiciuxf3n-de-completado-definition-of-done---dod}

La Definición de Completado describe formalmente cuándo un incremento cumple los criterios de calidad exigidos, es decir, el conjunto de verificaciones que una historia de usuario debe superar para declararse terminada y lista para su entrega o integración \citep{schwaber2020}. En este proyecto, una historia de usuario se considerará terminada si cumple las condiciones que enumera la Tabla~\ref{tab:definicion-de-completado-dod-para}:

\captionof{table}{Definición de Completado (DoD) para el proyecto}\label{tab:definicion-de-completado-dod-para}

\begin{longtable}[]{@{}
  >{\raggedright\arraybackslash}p{(\columnwidth - 2\tabcolsep) * \real{0.2903}}
  >{\raggedright\arraybackslash}p{(\columnwidth - 2\tabcolsep) * \real{0.7097}}@{}}
\toprule\noalign{}
\begin{minipage}[b]{\linewidth}\raggedright
\textbf{Criterio}
\end{minipage} & \begin{minipage}[b]{\linewidth}\raggedright
\textbf{Descripción}
\end{minipage} \\
\midrule\noalign{}
\endhead
\bottomrule\noalign{}
\endlastfoot
\textbf{Desarrollo completado} & La funcionalidad o actividad comprometida ha sido implementada.

Una historia relacionada al modelo predictivo implica que el modelo ha sido entrenado, evaluado mediante las métricas definidas en el proyecto, documentado y validado por el equipo.

Una historia asociada a la optimización de rutas deberá generar recorridos válidos y sus resultados deberán ser consistentes con los criterios de éxito establecidos en la fase de Definir. \\
\textbf{Validación realizada} & Se verificó que los resultados cumplen los criterios de aceptación definidos. \\
\textbf{Documentación actualizada} & Los cambios realizados se encuentran documentados dentro del proyecto. \\
\textbf{Evidencia disponible} & Existen pruebas, capturas, resultados o entregables que demuestran la ejecución de la actividad. \\
\textbf{Revisión del Product Owner} & El Product Owner valida que el resultado satisface la necesidad planteada. \\
\textbf{Integración exitosa} & El entregable puede utilizarse dentro de la solución sin generar conflictos con otros componentes. \\
\end{longtable}

\fuente{Fuente: Elaboración propia}

\subsection{Principios LEAN y métricas}\label{principios-lean-y-muxe9tricas}

La metodología LEAN se empleará como complemento de Scrum para validar que el producto mínimo viable (MVP) desarrollado genere valor real para la organización. Bajo este enfoque, la construcción del MVP estará orientada a comprobar hipótesis relacionadas con la mejora de la planificación logística y la reducción de ineficiencias operativas identificadas durante la fase de análisis. El ciclo de trabajo seguirá el principio Construir--Medir--Aprender, permitiendo evaluar continuamente los resultados obtenidos y realizar ajustes cuando sea necesario.

\subsubsection{Las etapas del proceso de desarrollo de software LEAN}\label{las-etapas-del-proceso-de-desarrollo-de-software-lean}

La aplicación eficaz de LEAN al desarrollo de software comprende cinco actividades \citep{paladiy2026}: identificar el valor que el producto ofrece al cliente; mapear el flujo de valor, distinguiendo qué partes del proceso o qué funciones aportan valor y cuáles no; crear flujo, de modo que el trabajo no se acumule entre etapas ni aparezcan cuellos de botella; establecer el sistema de tracción, para que las actividades del equipo las impulse lo que el cliente valora y no lo que pudiera hacer falta en el futuro; y buscar la perfección mediante la revisión y mejora continua del propio proceso. Como Scrum se apoya en esta misma filosofía, las cinco etapas se materializan en sus eventos, tal como se describe a continuación.

\begin{itemize}
\item
  \textbf{Identificar valor.} Las historias de usuario de la lista de tareas pendientes y el orden de prioridad asignado se originan en los requerimientos de los beneficiarios del sistema —el gerente propietario de Cafetalino, el operador de recarga y los clientes—, recogidos en las entrevistas y la encuesta, y transmitidos al equipo por Francisco Camacho (Product Owner).
\item
  \textbf{Mapear el flujo.} El análisis del flujo de trabajo y los ajustes necesarios para evitar implementar funciones innecesarias se realizan en las reuniones de Scrum y en las de planificación y revisión de cada iteración, estas últimas en el contexto más amplio del producto final.
\item
  \textbf{Crear flujo.} Las reuniones de Scrum permiten detectar y resolver cuellos de botella en el día a día, pero el evento decisivo es la retrospectiva de cada iteración, que analiza el desempeño del equipo y ajusta el proceso de desarrollo.
\item
  \textbf{Establecer el sistema de tracción.} Se manifiesta en las reuniones de planificación, donde la selección y priorización de historias responden a los requerimientos del gerente propietario y del operador de recargas, canalizados a través de las directrices del Product Owner.
\item
  \textbf{Buscar la perfección.} Ocurre en las reuniones de planificación y de revisión, orientadas a la calidad del sistema, y sobre todo en las retrospectivas, donde el desempeño del equipo y el proceso de desarrollo son objeto de análisis y mejora continua.
\end{itemize}

\subsubsection{Hipótesis de valor}\label{hipuxf3tesis-de-valor}

Con base en los hallazgos obtenidos durante las fases de Empatizar y Definir, se establecieron las siguientes hipótesis de validación:

\begin{enumerate}
\def\labelenumi{\arabic{enumi}.}
\item
  El uso de modelos predictivos basados en datos históricos permitirá reducir las visitas de reabastecimiento innecesarias realizadas por el personal operativo.
\item
  La planificación de rutas basada en la demanda proyectada permitirá disminuir las modificaciones extraordinarias de recorrido ocasionadas por eventos de desabastecimiento.
\item
  La aplicación de técnicas de aprendizaje automático proporcionará una capacidad de predicción superior al método actualmente utilizado por la organización, basado principalmente en experiencia operativa y programación periódica de visitas.
\item
  La utilización del sistema propuesto contribuirá a reducir incidencias relacionadas con desabastecimiento y mejorará la continuidad del servicio ofrecido a los usuarios finales.
\end{enumerate}

\subsubsection{Métricas}\label{muxe9tricas}

En el proceso de desarrollo de software LEAN, los indicadores de rendimiento ayudan a medir el valor, el tiempo y la eficiencia de los procesos, lo que permite evitar desperdicios \citep{paladiy2026}. Para este proyecto se aplicarán las siguientes métricas:

\begin{itemize}
\item
  Tiempo promedio que se tarda en completar una historia de usuario, desde que inicia su desarrollo hasta que es aprobada por el equipo.
\item
  Porcentaje de Historias de Usuario aprobadas en la revisión de la iteración, sin requerir ajustes, ni trabajo adicional.
\item
  Porcentaje de incrementos que provocan errores de regresión; es decir cambios que rompen funcionalidades de historias de usuario completadas en anteriores iteraciones.
\item
  Para validar las hipótesis planteadas y medir la calidad del sistema desarrollado se utilizarán indicadores cuantificables alineados con los criterios de éxito de la Tabla \ref{tab:criterios-de-exito}, establecidos en la etapa Definir.
\end{itemize}

\subsubsection{Estrategia Construir -- Medir -- Aprender}\label{estrategia-construir-medir-aprender.}

En el desarrollo del sistema propuesto, la aplicación de LEAN se realizará mediante la estrategia cíclica de validación construir-medir-aprender, que es el núcleo de esta metodología y busca transformar una idea en un MVP mediante ciclos rápidos de iteración, evaluando la respuesta del usuario y decidiendo si se debe reorientar la lógica de desarrollo del proyecto o continuar con la línea de trabajo aplicada. Sus tres etapas, en el marco del proyecto, se describen a continuación.

\textbf{Construir} el producto mínimo viable, constituido por el modelo predictivo de demanda, el componente de optimización de rutas y la interfaz de visualización para apoyo a la toma de decisiones. \textbf{Medir} su desempeño mediante los indicadores de éxito de la Tabla \ref{tab:criterios-de-exito}, definidos en la etapa Definir, comparando los resultados obtenidos contra la línea base que hoy utiliza Cafetalino. Y \textbf{aprender} de esos resultados para identificar oportunidades de mejora, ajustar los parámetros de los modelos y refinar la solución antes de una posible implementación operativa.

En conjunto, Scrum aporta la estructura con la que se organizará y verificará el trabajo —roles, artefactos, eventos y las definiciones de Listo y Completado—, mientras que LEAN aporta el criterio con el que se decidirá si lo construido genera valor para Cafetalino. Ambos marcos convergen en el mismo punto: cada incremento debe poder medirse contra la operación actual de la empresa antes de darse por bueno. El siguiente capítulo describe cómo se construyen los cuatro componentes del sistema bajo esta metodología, con qué tecnologías y con qué estimación de tiempo y coste.
